\section{Real valuations and a finite constrained optimization}
\label{c:optimization-section}

The section estimates are now in place. We seek a valuation whose
discrepancy is small, whose values on all sections satisfy a lower bound,
and whose selected initial coordinates remain independent. The first task
is to make this optimization finite and attain its minimum on the original
function field. The second is to extract a strict divisorial inequality
from optimality; that inequality is the input for Cartier trace in the
next section.

We return to the actual function field
\(K=k(X)=k(Z)\), the fractional form \(\theta\), the valuation \(v\),
and the nested multiplicative systems \(V_m\) constructed above.
All physical degrees in this and the following sections belong to
\(n\Z_{\ge0}\).  If \(0\ne F\in V_m\), its \emph{assigned cost} is
\[
                         b(F)=m+v(F).
\]
The assigned degree is part of this notation.  We have
\(V_0=k[Z]\) and \(\Frac(V_0)=K\).  The only zero-cost elements
are the degree-zero functions invertible at \(z\).  The preceding
weight estimates and \eqref{eq:A7} give
\[
              m\le Cb(F)\quad\text{when }b(F)>0,
\]
and a uniform positive lower bound for positive costs along the
chosen subsequence.  Constants may depend on the fixed subsequence
bounds already selected.

The earlier normalized-volume minimization supplied the graded model
underlying these section estimates.  In the present optimization,
\(v\) remains fixed and determines the costs \(b(F)\), while the
candidate valuation varies on the actual field \(K\).  The objective
balances its angular discrepancy against a common lower bound for
section values, with prescribed independent initial coordinates.

\subsection{Constants and full initial fields}

We may, and do, take \(k\) countable and algebraically closed of
characteristic zero.  Here is the descent needed for this choice.
Descend the countably many members of the sequence and all their finite
data simultaneously from \(\C\): the points \(z\), the sections \(h\),
the graded algebra presentations and generators with their lifts,
the monomial charts for \(v,\widetilde v\), the models and contractions,
the lattice and field identifications, the forms and divisors, and the
chosen probes.  Retain the numerical valuation weights.  These data
are defined over the algebraic closure of a countably generated
subfield of \(\C\), which is countable.

We explain why the assertions used below survive this descent.
On an independent-normal-weight chart as in \eqref{eq:A6}, the full
initial field is the stratum field with independent normal variables.
The leading-monomial congruences and their nonzero residue coefficients
are preserved on extending constants.  For any finite-dimensional
space of functions over \(k\), there are only finitely many values:
elements of distinct values are linearly independent.  Taking bases
on its successive value slices gives a basis with independent initials.
These initials remain independent after extension to \(\C\), since
the stratum is geometrically integral and its initial coordinates
extend injectively.  This computes the whole filtered space after
extension as well.  Sections of the reflexive powers defining the
graded affine algebra commute with constant extension by flat base
change; alternatively compute them on the smooth locus.  Taking
finite spans therefore proves that the associated graded algebras,
the prescribed lifts, and the field identifications extend as required.

In particular the leading-space dimensions, widths, and nonemptiness
assertions descend.  Retain \(S_1,H_1,A_1^+\), \eqref{eq:A9}, and the
terminal field \(K_f\).  A divisorial test over \(k\) extends from a
smooth divisor chart to \(\C\), with its orders and discrepancies
unchanged, so \eqref{eq:C4} continues to hold.  If its restriction to
\(K_f\) is trivial, the residue embedding of \(K_f\) stays injective
after tensoring and taking fractions, so this condition is preserved
too.  The \(\epsilon\)-lc condition on \(X\) and the lc condition
defined by \(h\) likewise descend.  A contradiction over \(k\)
will consequently suffice.  We shall enlarge constants again when
forming geometric generic fibres.

All valuations are trivial on their stated ground field.  For a real
Abhyankar valuation \(w\) of \(K\), including the zero valuation,
write \(K_w\) for the fraction field of its associated graded ring,
and \(\overline F\) for a homogeneous initial.  Thus \(K_w\) is the
\emph{full initial field}, not merely the residue field.  Write \(E\)
for the residue field of \(w\).  A valuation is \emph{Gaussian} on a
list if the initials of that list are algebraically independent over
the constants.  Equivalently, every nonzero polynomial in the list
has the termwise minimum value.

\subsection{The barrier and its finite test}

Fix one sequence member, positive real numbers
\(B_*,\kappa,q_*,q_{\rm bar}\) with \(q_{\rm bar}<q_*\), and an
algebraically independent list \(y=(y_j)\) of rational functions,
possibly empty.  The symbol \(B_*\) is a budget scalar, not a
boundary.  Define
\begin{equation}\label{eq:E3}
 q(w)=\min\left\{q_*,
    \inf_{\substack{m,\;0\ne F\in V_m\\ b(F)>0}}
                  \frac{w(F)+B_*m}{b(F)}\right\},
 \qquad
 \mathcal L(w)=A_\theta(w)-\kappa q(w).
\end{equation}
For a real parameter vector \(t\), minimize \(\mathcal L\) over
real quasi-monomial \(w\) with \(q(w)\ge q_{\rm bar}\), Gaussian
on \(y\), and \(w(y)=t\).  A feasible valuation is centred over
\(z\) on \(X\): its bound on \(V_0=k[Z]\) gives positivity on the
maximal ideal of \(z\), zero on functions invertible there, and
properness supplies the centre on \(X\).  The lc boundary in
\eqref{eq:A13} then gives \(A_\theta(w)\ge0\).

\begin{lemma}[Finite barrier test]\label{c:finite-barrier}
There are finitely many pairs \((m_i,F_i)\), with
\(b_i=m_i+v(F_i)>0\), such that whenever their minimum ratio,
including \(q_*\), is at least \(q_{\rm bar}\), it equals \(q(w)\).
An index attaining the minimum is called active; the cap is an
additional active entry when \(q(w)=q_*\).
\end{lemma}

\begin{proof}
Take the degree-\(n\) Veronese \(\mathcal A'\) of the actual section
algebra \(\mathcal A\).  Write \(o'\) for the image of \(o\) in
\(\Spec\mathcal A'\).  The initial algebra of \(\mathcal A'\)
at \(\widetilde v\) is the corresponding Veronese of \(R\):
\(\widetilde v\) is degree-torus invariant, and linear combinations
of torus translates separate the physical degree characters.
This algebra is finitely generated, with value-zero part consisting
only of constants.  Choose lifts
\(s_i=h^{m_i/n}F_i\) of homogeneous initial generators and include
ordinary homogeneous algebra generators of \(\mathcal A'\).
In degree zero subtract their value at the point \(o'\), so all
these generators vanish there and have positive cost.

Every homogeneous \(s=h^{m/n}F\) vanishing at \(o'\) has an exact
polynomial representative in these generators, all of whose terms
have physical degree \(m\) and cost at least
\(b=\widetilde v(s)\).  To prove exactness, subduct by the initial
generators, retaining physical degree, to arbitrarily high cost.
The residual value strictly increases; only finitely many costs
occur below a fixed bound, by positivity and finite graded
generation.  The fixed-germ log Izumi estimate at the bare klt
cone point \(o\) then puts a sufficiently high residual in any
prescribed power of its maximal ideal.  The powers are primary
also on the affine ring, so this assertion is global.  It may
equally be checked after extension to \(\C\) and contracted.

The same conclusion holds in \(\mathcal A'\).  Indeed the radical
of its point ideal extended to \(\mathcal A\) is the point ideal
of \(o\), as seen by taking degree powers; moreover the Veronese
inclusion splits as \(\mathcal A'\)-modules.  Thus a sufficiently
high residual belongs to any prescribed power of the point ideal
in \(\mathcal A'\).  Such a power has representatives by
sufficiently long monomials in the ordinary generators, and hence
by terms of high cost.  Project to physical degree \(m\).
This finishes subduction by an exact finite representative.

Divide by \(h^{m/n}\).  If the finite ratios are bounded below
by a positive number \(q\), every term satisfies
\(w(\text{term})+B_*m\ge q\,(\text{its cost})\ge qb(F)\).
Taking the minimum gives the same bound for \(F\).
Zero-cost units cause no additional condition.  This proves the
asserted finite test.
\end{proof}

\subsection{Monomial models and variations}

The finite test replaces the section inequalities by finitely many
actual functions.  We still need to control discrepancy while
retracting a valuation to a finite skeleton, and to construct variations
with controlled probe values and preserved Gaussianity.  The
following monomialization and differential arguments provide these
two kinds of control.

We will also use ordered-group valuations satisfying Abhyankar
equality.  The relative Abhyankar inequality says that, for an
extension of valued function fields, the increment of rational rank
plus the increment of residue transcendence degree is at most the
transcendence degree of the extension.  Indeed choose values
independent modulo the old values over \(\Q\), and value-zero
elements with residues independent over the old residue field.
Their monomials are independent, by first comparing values and then
residues.  It follows that Abhyankar equality passes by restriction
to subfunction-fields.  The transcendence degree of the full initial
field is rational rank plus residue transcendence degree: choose
homogeneous elements whose values rationally span the value group;
all remaining homogeneous elements become algebraic after taking
monomial ratios.

\begin{lemma}[Ordered monomialization]\label{c:ordered-monomialization}
An ordered-group valuation of a characteristic-zero function field
satisfying Abhyankar equality is monomial on a smooth projective
model, which may dominate a prescribed model and resolve prescribed
finite divisor data.  Its centre is the generic point of a transverse
stratum; the positive normal-parameter values are rationally independent.
Its full initial field is the stratum field with independent normal
variables.
\end{lemma}

\begin{proof}
Choose functions with rationally independent values spanning the
rational value group, and value-zero functions with residues forming
a complementary transcendence basis.  Resolve the divisors of the
first functions and the maps to \(\PP^1\) of the second functions,
together with the prescribed finite data.  At the centre, the number
of crossing components is at least the rational rank: the chosen
functions are units times Laurent monomials in their local equations.
The closure of the centre has dimension at least the residue
transcendence degree, because the chosen residue functions are regular
and invertible there.  Equality of their sum with the field dimension
forces the centre to be generic in a transverse stratum and the normal
values to be rationally independent.

Complete the regular local ring at that generic point and expand over
a coefficient field.  The support of a nonzero series has finitely
many componentwise-minimal exponents.  Since all normal weights are
positive, its least weight is attained among them; rational
independence makes the least exponent unique.  This also works in the
ordered group.  The ideal of monomials of strictly higher weight is
finitely generated.  Faithful flatness contracts the congruence modulo
this ideal to the local ring.  Hence the function equals its least
monomial times a residue-coefficient lift, modulo strictly higher
weight.  This proves the value formula and the initial-field
description.  Independence of the normal values also identifies the
degree-zero field with the stratum field.
\end{proof}

For a real \(w\), Lemma~\ref{c:ordered-monomialization} gives
\[
                             K_w=E(\mathbf z),
\]
where the homogeneous Laurent variables \(\mathbf z\) have degrees
forming a basis of the value lattice.  This lattice is free of finite
rank by the same chart.  Any homogeneous elements of those basis
degrees can serve as the variables: correct characters to degree zero
and use independence of distinct degrees.  Fractional forms are
interpreted by taking their \(n\)-th powers and dividing form orders
by \(n\).  In particular \(\theta_w^n\) is homogeneous for this value
torus, by the leading-form rule \eqref{eq:A6}.

\begin{lemma}[Logarithmic determinant]\label{c:log-determinant}
If \(a_1,\ldots,a_d\) are a transcendence basis and \(w\) is a real
Abhyankar valuation in characteristic zero, then
\[
                  A_{\bigwedge_i d\log a_i}(w)\ge0.
\]
Equality holds exactly when the initials of the \(a_i\) are
algebraically independent.
\end{lemma}

\begin{proof}
Include their divisors in the monomialization.  Use logarithmic
differentials of the normal parameters and of unit lifts of a
separating basis on the stratum.  Their wedge has form discrepancy
zero.  Each \(a_i\) is a unit times a normal Laurent monomial.
The logarithmic differential matrix has entries of nonnegative
value, and its reduction is the corresponding matrix of initial
logarithmic differentials.  Logarithmic normal derivatives of a unit
reduce to zero; tangential derivatives differentiate the residue
coefficient.  The determinant therefore has value zero exactly
when the initial differentials are independent, equivalently when
the initials are algebraically independent in characteristic zero.
\end{proof}

\begin{lemma}[Invariantization]\label{c:invariantization}
Let \(Q^*\) be a real Abhyankar valuation of \(K_w=E(\mathbf z)\).
Put \(P=Q^*|_E\), and let \(Q\) be the Gauss valuation with restriction
\(P\) and the same values as \(Q^*\) on \(\mathbf z\).  Then \(Q\)
is Abhyankar, agrees with \(Q^*\) on every homogeneous element, and
\begin{equation}\label{eq:E1}
                    A_{\theta_w}(Q^*)\ge A_{\theta_w}(Q).
\end{equation}
Equality implies \(Q^*=Q\).  Gaussianity of \(Q^*\) on a list of
homogeneous probes is preserved by \(Q\).
\end{lemma}

\begin{proof}
The restriction \(P\) is Abhyankar by the relative inequality.
The Gauss extension is Abhyankar because its variable initials
are independent over the full field \(E_P\).  Choose functions
\(x\) in \(E\) with full independent initials at \(P\).
The coefficient of \(\theta_w^n\) in the \(n\)-th power of
\(d\log x\wedge d\log\mathbf z\) is homogeneous, so its value is
the same for \(Q\) and \(Q^*\).  The logarithmic differential
contribution is zero at \(Q\), and is nonnegative at \(Q^*\) by
Lemma~\ref{c:log-determinant}.  Equality requires independent
initials of \(x,\mathbf z\).  As \(E_P\) is algebraic over the
field generated by the \(x\)-initials, the \(\mathbf z\)-initials
are then independent even over \(E_P\).  This gives the termwise
Gauss rule at \(Q^*\), proving equality of valuations.

For the last assertion, group a polynomial in the probes by its
torus characters.  The two valuations agree on each homogeneous
part, and \(Q\) takes their minimum.  Gaussianity at \(Q^*\)
therefore gives the same polynomial minimum at \(Q\).
\end{proof}

\begin{lemma}[Flattening a secondary valuation]\label{c:flattening}
For the invariant valuation \(Q\) of
Lemma~\ref{c:invariantization}, there are real monomial valuations
\(w_\lambda\), for small \(\lambda>0\), such that
\begin{equation}\label{eq:E2}
 \begin{split}
 w_\lambda(F)&=w(F)+\lambda Q(\overline F),\\
 A_\theta(w_\lambda)&=A_\theta(w)+\lambda A_{\theta_w}(Q).
 \end{split}
\end{equation}
For each fixed nonzero rational function the first identity holds
for all sufficiently small \(\lambda\), with the interval allowed
to depend on that function.  If \(w\) is Gaussian on a prescribed
finite list \(y\), and \(Q\) is Gaussian on \(\overline y\), then
\(w_\lambda\) may be kept Gaussian on \(y\).
\end{lemma}

\begin{proof}
Give \(F\) the lexicographically ordered value
\((w(F),Q(\overline F))\).  Grouping tied first values proves the
valuation inequality.  Its initials identify with double initials:
on each first homogeneous slice the second comparison is exactly
the definition, while distinct first degrees remain separated by
their \(\mathbf z\)-characters, with coefficients in \(E_P\).
Multiplication is preserved.  Lifts to \(K\) of the variables
\(\mathbf z\) and of the functions \(x\) used above have full
independent lex initials.  Such lifts exist for all homogeneous
elements, including degree-zero residue elements.  The lex valuation
therefore satisfies Abhyankar equality.

Monomialize it by Lemma~\ref{c:ordered-monomialization}.  Replace each
normal weight by its first coordinate plus \(\lambda\) times its
second coordinate.  These weights are positive for small
\(\lambda>0\).  For any fixed numerator and denominator, the finitely
many componentwise-minimal support exponents show that its former
least exponent remains least eventually.  The resulting initial
with its stratum coefficient and normal variables is unchanged.
The normal initial variables remain independent even if the new
real values have ties, since a monomial valuation retains all
least-weight terms.  This proves the value assertion and preserves
independence of the lifted full probes.  Double initials likewise
preserve independence of \(y\) under the two stated Gaussian
conditions.

Express \(\theta^n\) in the logarithmic basis of those full probes.
Its coefficient computes the discrepancy at \(w\) by
\eqref{eq:A6}, and its initial is the coefficient of \(\theta_w^n\).
The logarithmic basis contributes zero at \(w_\lambda\) and at \(Q\).
Apply the value identity also to this one fixed coefficient and
divide by \(n\); the second identity follows.
\end{proof}

\subsection{Retraction and compact optimization}

We now apply these valuation tools to the finite barrier data.
Retraction will restrict the valuations to finitely many skeleton
cones; the comparison with \(A_X\) in \eqref{eq:A13} will exclude
escape to infinity on an objective sublevel.

\begin{lemma}[Retraction and strict discrepancy drop]
\label{c:skeleton}
Resolve on a smooth projective model the divisor supports of
\(\theta^n,F_i,y_j\).  Retraction to its normal crossing skeleton
preserves the finite barrier data, probe values and probe
Gaussianity, and does not increase \(A_\theta\).  The discrepancy
decreases strictly for a real Abhyankar valuation outside the
skeleton.  On every closed skeleton cone, each fixed rational
function has a continuous finite piecewise-linear value, with
a finite Lipschitz constant.
\end{lemma}

\begin{proof}
At a centre, keep the weights on the crossing normal equations
and move to the generic stratum; if there are no such equations,
the retraction is zero.  Retraction can only decrease the value
of a regular function.  To see this with the required control at
faces, let \(x\) be the normal parameters.  A monomial ideal of
positive threshold contracts from the generic stratum to the
local ring.  In its completed quotient there is a finite free
module over the formal series in the other parameters.  A
function outside the stratum prime acts injectively: its reduction
in the other parameters is nonzero, and the remaining normal part
acts nilpotently.  The monomial ideal is therefore primary to
the stratum prime, and faithful flatness gives contraction.
Membership in this ideal characterizes the threshold at the
generic stratum and bounds the value at the original centre.

On an open cone, normal expansion computes values by the minima
over finitely many componentwise-minimal support exponents of
regular numerators and denominators.  At a face apply the same
primary-ideal argument to just the remaining positive-weight
parameters.  It proves that these formulas extend to the face;
at the origin both values are zero.  This gives continuity, finite
piecewise linearity and a Lipschitz bound.  Divisors included in
the resolution have linear values on each cone.  So does
\(A_\theta\), using the local logarithmic top generator.

The retraction preserves all resolved divisor values.  For a
polynomial in \(y\), clear poles by a normal monomial, since each
\(y_j\) is a unit times a normal Laurent monomial.  Its retracted
value is at most its old value and at least its termwise minimum.
If \(w\) is Gaussian these agree, proving preservation.
Within one open cone the initial formulas of \(y\) in stratum
coefficients and independent normal variables do not depend on
the positive weights.  Gaussianity thus selects whole open cones.
It persists on their faces by continuity of all polynomial values.

For the discrepancy comparison, choose supplementary regular
functions \(z'\) with étale differentials at the centre, including
a separating residue basis, and put
\(\lambda_0=d\log x\wedge dz'\).  The coefficient of \(\theta^n\)
relative to \(\lambda_0^n\) is a unit-monomial.  The difference
from the retracted discrepancy is \(A_{\lambda_0}(w)\).
Switching the regular supplementary factors to logarithmic ones
and applying Lemma~\ref{c:log-determinant} makes this nonnegative.
If the centre is proper in its stratum, a supplementary parameter
vanishing there has positive value, making the difference strict.

If the centre is generic in the stratum, use unit lifts of a
separating stratum basis.  Equality in the determinant test makes
those initials and the normal initials independent.  The normal
initials are independent even over the embedded stratum field,
which is algebraic over that basis field.  Least normal terms
therefore cannot cancel; contracting their congruence modulo
strictly higher monomial weight proves that \(w\) itself is the
monomial retraction.  This proves strictness outside the skeleton.
The proof is unchanged over other algebraically closed
characteristic-zero constant fields.
\end{proof}

\begin{proposition}[Compact optimization and continuation]
\label{c:optimization}
The optimization in \eqref{eq:E3} is a finite polyhedral
optimization with compact sublevels, and attains its minimum
whenever feasible.  Every optimizer lies on the fixed skeleton.
At any parameter admitting an optimizer with \(q>q_{\rm bar}\),
the minimum is locally finite and continuous.  On full-dimensional
generic parameter cells it is affine.  Through every optimizer at
such a generic parameter there is a local affine continuation by
optimizers within its same open stratum cone and with the same
active mask.
\end{proposition}

\begin{proof}
By Lemmas~\ref{c:finite-barrier} and \ref{c:skeleton} we may optimize
on a finite union of closed polyhedral pieces, with affine
projection to \(t\) and piecewise affine objective.  Strict
retraction excludes every optimizer outside them.

Suppose a sublevel, allowing \(t\) to vary, were unbounded.
Normalize an escaping sequence in one closed cone by the norms
of its weights and take a limit of norm one, defining a nonzero
monomial valuation \(w'\).  The barrier and continuity give
\(w'(F)\ge0\) for every \(0\ne F\in V_m\), including degree-zero
units.  The bounded cap in \eqref{eq:E3} gives
\(A_\theta(w')\le0\).  Nonnegativity on \(k[Z]\), and properness
over the affine base, give a centre on \(X\).
Equation~\eqref{eq:A13}, valid at such centres by freeness of
sufficiently divisible multiples over the affine base, implies
\(A_X(w')\le0\).  This is impossible: on a monomializing resolution
a nonzero quasi-monomial valuation centred on a klt variety has
strictly positive discrepancy.  The sublevels are compact, proving
attainment.

For continuity, let \(w\) be an optimizer with strict barrier
slack.  Give \(k(\overline y)\) arbitrary rational-direction
Gauss weights.  Extend this valuation to \(K_w\) by completing
to a transcendence basis with freely weighted variables and then
extending over the finite algebraic extension.  The relative
Abhyankar inequality gives equality, and the extended value
group is contained in the rational span of the starting real
weights, so the extension remains real.  Invariantize by
Lemma~\ref{c:invariantization} and flatten by
Lemma~\ref{c:flattening}.  The finite test makes a small segment
in the prescribed direction feasible, with cost tending to the
optimum.  Retraction can only improve it.

A compact skeleton sublevel cut strictly above that optimum
projects to a finite union of closed polytopes.  The tangent
cones at the parameter contain all rational directions and hence,
by closedness and density, all directions.  Near the point only
the active inequalities of the incident polytopes matter.
Its projection therefore contains a neighbourhood.  Applying
this for sublevels arbitrarily close above the optimum proves
local finiteness and upper continuity.  Compactness proves lower
continuity.

Finally subdivide a compact sublevel into finitely many polyhedra
with affine cost and fixed active mask and stratum on each
relative interior.  Projection to parameter and cost gives the
piecewise-affine minimum.  Avoid its walls and all
lower-dimensional parameter projections of domain faces.
The affine hull of the relative interior containing any
optimizer then projects with full rank.  Choose a local affine
right inverse through that optimizer, staying in the same
relative interior.  Its cost minus the optimal value is affine,
nonnegative nearby and zero centrally, hence identically zero.
This is the required continuation.  In the applications below,
strict objective bounds maintain strict barrier slack; compactness
and continuity then also control limits along nongeneric segments.
\end{proof}

\subsection{A strict budget on a geometric generic fibre}

Compactness and continuation allow us to choose and vary optimizers.
We next examine a second valuation of their full initial field. If it
preserves the old initial coordinates, flattening gives an admissible
variation on the original field. Optimality controls its discrepancy;
a tie rule will make the resulting inequality strict after adding
finitely many small auxiliary terms.

Fix parameters admitting an optimizer with \(q>q_{\rm bar}\),
and let \(\mathcal M\) be the compact set of all optimizers.
We first maximize \(q\) on \(\mathcal M\).  The finite barrier
test makes \(q\) continuous there, so its maximum \(q_{\max}\)
is attained and satisfies \(q_{\max}>q_{\rm bar}\).  The locus
\[
             \mathcal M_{\max}=\{u\in\mathcal M:q(u)=q_{\max}\}
\]
is therefore nonempty and compact.
Choose the coefficients \(c_F>0\), indexed by the countable nonzero
angular functions, so rapidly decreasing that
\[
                          \Psi(u)=\sum_F c_Fu(F)
\]
converges absolutely and uniformly on \(\mathcal M\), and that
the weighted Lipschitz constants are summable on each of the
finitely many closed cones.  This is possible by
Lemma~\ref{c:skeleton}.
We then choose \(w\in\mathcal M_{\max}\) maximizing \(\Psi\)
there.  In particular \(q(w)=q_{\max}>q_{\rm bar}\).  We call
this an optimizer with the stated tie rule.

Take independent homogeneous probes
\(g=(\overline y,g_{\rm add})\) in \(K_w\).  The optional
\(g_{\rm add}\) are not optimization constraints.  Let \(L\)
be any geometric generic fibre field over \(k(g)\): choose a
component of the compositum with
\(k'=\overline{k(g)}\) and work over \(k'\).  Divide \(\theta_w\)
by the ordinary base wedge \(\bigwedge_j dg_j\), interpreting
fractional forms through powers, to obtain
\(\theta_{w,\mathrm{rel}}\).

\begin{proposition}[Relative optimality and its equality case]
\label{c:relative-budget}
For every nontrivial real Abhyankar valuation \(Q^\dagger\) of
\(L/k'\),
\begin{equation}\label{eq:E4}
 A_{\theta_{w,\mathrm{rel}}}(Q^\dagger)
       \ge \kappa\min_{i\ \mathrm{active}}
                       \frac{Q^\dagger(\overline F_i)}{b_i}.
\end{equation}
An active cap contributes the entry zero.  If equality holds,
then \(Q^\dagger(\overline F)<0\) for some actual angular
function \(F\).
\end{proposition}

\begin{proof}
Restrict \(Q^\dagger\) to a valuation \(Q^*\) of \(K_w\).
The extension \(L/K_w\) is algebraic, and its enlarged constants
contribute \(|g|\) to residue transcendence degree over \(k\).
The relative inequality thus shows that \(Q^*\) has Abhyankar
equality.  The \(g_j\) have value zero and independent initials.
Complete them, using functions of \(K_w\), to a list with full
independent initials.  Their supplementary initials remain
independent over \(k'\) at \(Q^\dagger\): the graded-field
embedding is injective, and the constant extension is algebraic
over \(k(g)\).  Expressing the form in \(dg\) wedged with
logarithmic supplementary differentials and using \eqref{eq:A6}
gives
\[
             A_{\theta_{w,\mathrm{rel}}}(Q^\dagger)
                              =A_{\theta_w}(Q^*).
\]
The coefficient is unchanged on dividing by the base wedge;
changing \(dg_j\) to \(d\log g_j\) contributes only value-zero
factors.

Invariantize \(Q^*\) to \(Q\).  It agrees on the homogeneous
active entries and preserves zero-value Gaussianity on the old
probes \(\overline y\).  Equation~\eqref{eq:E2} gives a path
\(w_\lambda\) preserving the old Gaussian parameters.  No
preservation of additional constraints is being required.  Let \(\mu\) be the
active minimum in \eqref{eq:E4}, including the zero entry if the
cap is active.  For small \(\lambda>0\), the finitely many test
ratios and the cap have minimum \(q(w)+\lambda\mu\).
Since \(q(w)>q_{\rm bar}\), this stays above the barrier;
Lemma~\ref{c:finite-barrier} identifies it with \(q(w_\lambda)\)
and makes the path feasible for the original optimization.
Original objective optimality and \eqref{eq:E1} prove the inequality.

Equality forces equality in invariantization, hence \(Q^*=Q\),
and an exact path of original optimizers \(w_\lambda\).
Suppose every angular initial had nonnegative \(Q^\dagger\)-value.
Then \(\mu\ge0\), including the possible cap entry.
Since the path consists of original optimizers and
\(q(w)=q_{\max}\), maximality of \(q\) forces \(\mu=0\).
Thus this equality path lies in \(\mathcal M_{\max}\).
Some angular value is positive:
otherwise all homogeneous initials, obtained as ratios of initials
from \(V_0\), have value zero, and the invariant Gauss rule makes
\(Q\) trivial.  Algebraicity would then make \(Q^\dagger\)
trivial as well.

This positive value contradicts the tie rule by differentiation
of \(\Psi\).  The interchange of sum and derivative is justified
as follows.  Strict retraction puts every small \(w_\lambda\)
on the fixed test skeleton.  Monomialize the lex valuation on a
model dominating that skeleton.  The flattened positive normal
weights have fixed centre for small positive \(\lambda\), so
their images use one skeleton cone.  Values of its finitely
many normal equations are affine in \(\lambda\) on a common
interval by \eqref{eq:E2}.  Their face limit is \(w\), by the
pointwise identities and closed-cone continuity.  Thus the
coordinate distance is \(O(\lambda)\).  The summable weighted
Lipschitz constants dominate all difference quotients; termwise
differentiation follows even though the eventual interval in
\eqref{eq:E2} depends on \(F\).  Every derivative is nonnegative
and at least one is positive, contradicting maximality of
\(\Psi\) on \(\mathcal M_{\max}\).
\end{proof}

\begin{corollary}[Strict rational budget]\label{c:strict-budget}
On the same geometric generic fibre, one can choose positive
rational numbers \(\alpha_i\) arbitrarily close to
\(\kappa/b_i\), a finite list of additional actual angular
functions \(U_j\), and arbitrarily small positive rational
\(e_j\), so that every prime divisor \(P\) over a proper model
satisfies
\begin{equation}\label{eq:E5}
 A_{\theta_{w,\mathrm{rel}}}(P)>
       \min_{i\ \mathrm{active}}\alpha_i P(\overline F_i)
            +\sum_j e_j\min\{0,P(\overline U_j)\}.
\end{equation}
The same zero-entry convention applies to an active cap.
The list may include field generators over \(k'\).
The added \(e_j\)-weighted terms and the errors from approximating
\(\kappa/b_i\) may meet any prescribed finite degree, value and
Lipschitz tolerances along a specified local affine optimizer
continuation.
\end{corollary}

\begin{proof}
Resolve the finite form and active-divisor data on a proper
relative model.  By strict retraction, the nonzero equality
directions in \eqref{eq:E4} lie on its normal crossing skeleton:
retraction preserves the minimum term and strictly lowers
discrepancy off the skeleton.  A nonzero valuation retracted
to zero has strictly positive discrepancy gap and is not an
equality direction.  Normalize the sum of the nonnegative
normal weights to one.  The equality locus is closed in the
finite compact union of these simplices.

Proposition~\ref{c:relative-budget} supplies an open cover of
this locus by conditions \(P(\overline U)<0\).  Choose a finite
subcover.  Add field generators: initials from \(V_0\) generate
\(K_w\), since \(K=\Frac(V_0)\).  Adding the indicated
two-element minimum systems with any positive small coefficients
makes the inequality strict everywhere, first with coefficients
\(\kappa/b_i\).

When a local affine continuation is specified, fix it before
choosing these coefficients.  It is a continuation by optimizers
of the original problem through \(w\); the tie rule is imposed
only at \(w\).  For each added or active actual function \(F\),
fix its assigned degree.  On a sufficiently small compact
parameter neighbourhood, its absolute values and Lipschitz
constant are bounded.  Let \(N(F)\) bound the sum of these
quantities and its assigned degree.  The added terms meet any
prescribed positive tolerance \(\eta\) by choosing the \(e_j\)
so small that \(\sum_j e_jN(U_j)<\eta\).  With no continuation
specified, retain just the degree and value bounds needed.

Resolve also the added functions.  The strict gap has positive
minimum on the normalized enlarged skeleton.  Each of the
finitely many active evaluations is bounded there, and
\[
       |\min_i a_i-\min_i b_i|\le\max_i|a_i-b_i|.
\]
Consequently a single sufficiently close rational perturbation
of the \(\kappa/b_i\) preserves strictness, including where
the minimizing index switches.  The cap remains a fixed zero
entry.  Retraction deals with every remaining direction,
including zero retraction.  Having chosen the \(e_j\) first,
make the rational perturbation still closer so that
\(\sum_i|\alpha_i-\kappa/b_i|N(F_i)<\eta\) and the resulting
strict margin is retained.  Finitely many prescribed tolerances
can be met simultaneously by taking \(\eta\) sufficiently small.

After clearing denominators, realize each term
\(e_j\min\{0,P(\overline U_j)\}\) by product powers of the
two-element system \(\{1,\overline U_j\}\).
The ratio \(\overline U_j\), rather than only one of its powers,
remains among ratios of the combined system: keep all other
factors fixed and replace one \(1\) by \(\overline U_j\).
This also supplies its asserted birationality.
\end{proof}
